As explored in \autoref{sec:v1-trade-memory}, the AT25xF256 was the better suited Renesas QSPI chip. The only difference between the S (single) and Q (quad) models the status of the QE bit on startup (whether QSPI is enabled by default). As QSPI will be used by default, the Q model was selected. OF the four potential sub-models, the \href{https://au.mouser.com/ProductDetail/Renesas-Dialog/AT25QF2561C-MWUB-T?qs=olJun0bQHM%2FqfVfqXrnLBg%3D%3D}{AT25QF2561C-MWUB-T} was chosen as it was the smallest (most compact) package in-stock. Future designs should use the \href{https://au.mouser.com/ProductDetail/Renesas-Dialog/AT25QF2561C-MUB-T?qs=olJun0bQHM%2Fuu3grdD8Csw%3D%3D}{AT25QF2561C-MUB-T} if it is stocked as it has a smaller package.
\nl
Section 2 (Pin Descriptions) of the \href{https://www.renesas.com/en/document/dst/at25sf2561cat25qf2561c-datasheet?r=25559126}{datasheet} was used to generate the pinout. As explored in \autoref{sec:v0.1-memory-remark}, the external pull up resistor on IO2 was not necessary in quad mode, but a pull up was required on NCS as there was not internal pull up in the chip.
\nl
As pointed out by \href{https://community.st.com/t5/stm32-mcus-products/is-internal-pull-up-the-same-as-external-pull-up-resistor/td-p/70477}{DTomi.1}, the rise/fall time of the signal is an important consideration for data lines. In Version 0.1, the average trace length was 21 mm, which corresponds to 24 pF at a trace width of 0.152 mm and conductor height of 0.1 mm (using the Saturn PCB Toolkit), with 6 pF added to account for input capacitance. Using \href{https://www.digikey.com/en/resources/conversion-calculators/conversion-calculator-time-constant}{a time constant calculator} we can see the time to fall to 0.7 VCC is 500 ns, far more than the 7.5 ns period of the QSPI chip operating at 133 MHz. If the fall time was too large, then I would risk the data being sent before the CS line was pulled down. To avoid this, a high value external pull up resistor (10k $\Omega$) was placed on NCS to minimise current draw..
\nl
The AT25QF2561C does not fully operate within the space temperature range of -40°C to 120°C so thermal management strategies were required to ensure the junction temperature does not exceed its specified maximum (85°C). As mentioned in \href{https://www.renesas.com/en/document/apn/thermal-resistance-and-its-applications}{this thermal resistance application note}, the best mitigation strategy is to utilise the thermal pad to draw heat away from the component. The datasheet does not specify a thermal resistance, so $\theta_{JC}=6.9$ C/W is assumed from the thermal characteristics of \href{https://www.onsemi.com/pdf/datasheet/ncp189-d.pdf}{this LDO} which shares a similar thermal pad. For an arbitrary ambient temperature of $T_A=80$°C and maximum $T_J=85$°C, then using the equation in case study 1 of the application note, the maximum current must be less than 700 mA. According to the datasheet, the maximum current sourced during any operation is 15 mA, so should be expected that the internal die temperature does not exceed 85°C through operation alone (and not through external heating), meeting the system's thermal requirement in \autoref{sec:v1-requirements}. Unfortunately, if the component is externally heated (from the Sun) then there is no alternative active cooling method to keep the junction temperature below 85°C and the component risks thermal damage.